\section{任务分析与总体思路}
情感信息分布在语言语境、声音韵律与视觉表现中。三类观测具有不同采样尺度，局部缺失和内容不一致又使简单拼接难以同时满足预测与解释要求。本文依据三项任务，分别建立原始素材的时间组织方法、局部缺失下的预测模型，以及具有局部贡献分解的融合模型。

问题一的关键在于保留观测并建立粒度和证据状态有区分的对应关系。原始音视频时钟提供物理联系，语音识别与声学边界提供文字对应候选，弱读连读情况下使用词组范围保留整体支持。特征输出同时包含真实长度、位置与有效性，供后续按样本或时间读取。

问题二研究给定特征中局部连续区间不可用时的预测行为。缺失在编码前施加，模型只汇总可见信息，并将可用比例送入融合层。类别平衡以改善三类性能分布为目标；模型效果通过缺失模态、比例、位置及形状的配对实验说明。

问题三将解释结构直接写入预测模型：单模态项和成对项构成输出，可靠性系数调整当前样本的融合尺度。贡献分数核算分类间隔或回归预激活，关键项保留原特征位置；复用问题一的定位方法关联媒体候选范围。

\begin{figure}[H]
\centering
\begin{tikzpicture}[font=\normalsize,box/.style={draw,rounded corners,align=center,text width=4.3cm,minimum height=1.25cm},>=Stealth]
\node[box] (a) at (0,0) {附件一原始素材\\文本、音轨、画面};
\node[box] (b) at (8,0) {问题一\\特征、掩码与时序关系};
\node[box] (c) at (0,-2.2) {附件二对齐特征\\训练与验证};
\node[box] (d) at (8,-2.2) {问题二\\局部缺失预测\\附件三30条结果};
\node[box] (e) at (8,-4.4) {问题三\\预测与贡献\\附件四20条结果};
\draw[->] (a)--(b);\draw[->] (c)--(d);\draw[->] (c.south) |- (e.west);
\draw[->,dashed] (b.east)--++(1,0) |- node[pos=.7,right,align=left]{媒体\\定位} (e.east);
\end{tikzpicture}
\caption{三项任务的实际输入与输出。虚线表示定位方法复用。}
\end{figure}

问题二、三使用题目提供的预计算特征训练，问题一新提取的矩阵作为独立特征结果保存。三个部分共享观测有效性与位置记录的处理思想，具体特征维度、标准化方法和预测模型分别定义。全文结果整理及构建代码使用OpenAI Codex（GPT-6）辅助；使用日期为2026年9月26日，版本发布日期未在本次环境提供。数值来自本地实验记录，公式与处理逻辑按相应实现核对。
